% AUTORIDAD DE FUENTE:
% expose_VIII_body.tex, líneas 2004--2029; índices de escaneo 300--302;
% folios de fuente 289--291. La corrección material de la capa francesa
% se conserva reversiblemente mediante \sisi.

\medskip
\noindent 6.2.1. Si $\overline G$ corresponde al par $(F,s)$, la
imagen inversa $f_n^*(\overline G)$ en $\overline G$ de la sección
$f_n$ de $\mathbb Z_{S,T}$ definida por la sección constante de
$\mathbb Z_S$ de valor $n\in\mathbb Z$ se identifica canónicamente,
como torsor bajo $G$, con la potencia $n$-ésima $F^{(n)}$ del
$G$-torsor $F$. El isomorfismo canónico de $G_U$-torsors
\[
 f_n^*(\overline G)|U\simeq G_U,
\]
definido por la primera relación de (5.7.1), viene determinado por la
sección $s^{(n)}$ de $F_U^{(n)}$, potencia $n$-ésima de la sección
dada $s$ de $F_U$. Esto explicita la descripción siguiente de las
secciones de $\overline G$ sobre $S$. Una sección $\overline g$ de
$\overline G$ sobre $S$ queda determinada por:

\smallskip
\noindent a) una sección $\underline n$ de $\mathbb Z_{S,T}$ sobre
$S$, que puede interpretarse como una sección de $\mathbb Z_T$, es
decir, como una función localmente constante sobre $T$ con valores en
$\mathbb Z$, o equivalentemente como una partición
\[
 T=\coprod_{i\in\mathbb Z}T_i
\]
de $T$ en una suma disjunta de abiertos indexados por
$i\in\mathbb Z$;

\smallskip
\noindent b) una familia de secciones
\[
 \begin{gathered}
  \overline g_i\in\Gamma(F^{(i)}/S_i'),\qquad i\in\mathbb Z,\\
  S_i'=S\setminus T_i',\qquad
  T_i'=\displaystyle\bigcup_{\substack{j\in\mathbb Z\\j\ne i}}T_j,
 \end{gathered}
\]
sujeta a la condición de compatibilidad siguiente. Para
$i,j\in\mathbb Z$, $\overline g_i$ y $\overline g_j$ coinciden sobre
$S_i'\cap S_j'=U$ después de identificar $F^{(i)}|U$ y $F^{(j)}|U$
con $G_U$ mediante las secciones $s^{(i)}$ y $s^{(j)}$,
respectivamente. Dicho de otro modo, debe existir una sección,
necesariamente única, $h_{ij}\in\Gamma(G/U)$ tal que
\[
 \overline g_i=h_{ij}s^{(i)},\qquad
 \overline g_j=h_{ij}s^{(j)}.
\]
Puesto que la descripción de \sisi{3.2}{6.2} de una extensión
$\overline G$ de $\mathbb Z_{S,T}$ por $G$ es manifiestamente
compatible con todo cambio de base $S'\to S$, lo anterior da una
descripción completa del haz $\overline G$ como funtor

%%%%%%%%%%  scan idx 301  |  folio 290  %%%%%%%%%%

del argumento $S'\in\operatorname{Ob}(\mathrm{Sch}_{/S})$.

\medskip
\noindent 6.2.2. Si $T$ es un esquema que es suma de esquemas conexos
no vacíos $T_i$, $i\in I$, la descripción de 6.2.1 de las secciones de
$\overline G$ sobre $S$ puede simplificarse. Tal sección queda
determinada por una familia de enteros $(n_i)_{i\in I}$ indexada por
$I$ y una familia de secciones
\[
 \overline g_i\in\Gamma(F^{(n_i)}/S_i'),\qquad i\in I,
\]
donde
\[
 S_i'=S\setminus T_i',
 \qquad
 T_i'=\coprod_{\substack{j\in I\\j\ne i}}T_j.
\]
Las $\overline g_i$ satisfacen la condición de compatibilidad
siguiente: para todo par $i,j\in I$, las restricciones de
$\overline g_i$ y $\overline g_j$ a $S_i'\cap S_j'=U$ coinciden bajo
las identificaciones de $F^{(n_i)}|U$ y $F^{(n_j)}|U$ deducidas de las
respectivas secciones $s^{(n_i)}$ y $s^{(n_j)}$ de estos torsors.

\medskip
\noindent 6.2.3. Finalmente, si se supone que $T$ es no vacío y
conexo, dar una sección $\overline g$ de $\overline G$ sobre $S$
equivale a dar un entero $n$ y una sección $\overline g$ de $F^{(n)}$
sobre $S$. Ya no hay ninguna condición de compatibilidad que imponer,
de modo que la sección $s$ de $F|U$ no interviene en la descripción del
grupo de secciones de $\overline G$ sobre $S$. (Naturalmente, $s$
vuelve a aparecer en cuanto se intenta describir los puntos de
$\overline G$ con valores en un $S$-esquema general $S'$.)

\medskip
\noindent 6.2.4. Como $\overline G$ es la unión de los
``subesquemas abiertos'' $f_n^*(\overline G)$, $n\in\mathbb Z$,
descritos en 6.2.1, se sigue que $\overline G$ es representable si y
solo si lo son los torsors $F^{(n)}$. Por ejemplo, basta con que $G$
sea representable y afín sobre $S$, por la teoría del descenso de
SGA 1 VIII 2.1.

%%%%%%%%%%  scan idx 302  |  folio 291  %%%%%%%%%%
